\documentclass[10pt,a4paper]{amsart}
\usepackage[margin=19mm]{geometry}
\usepackage{amsmath,amssymb,mathtools}
\usepackage[scr=rsfs,cal=euler]{mathalfa}
\usepackage{xfrac}
\usepackage[unicode,hidelinks,bookmarks=false]{hyperref}
\newcommand{\ie}{i.e.\ }
\newcommand{\cf}{cf.\ }
\newcommand{\resp}{respectively\ }
\newcommand{\Subsection}{\subsection}
\providecommand{\nobreakdash}{\nobreak\hbox{-}}
\setlength{\emergencystretch}{2em}
\multlinegap0pt
\usepackage{bm}
\usepackage{bbm}
\usepackage{dsfont}
\usepackage{xspace}
\usepackage{cancel}
\usepackage{mathtools}
\usepackage{amsthm}
\makeatletter
\@ifundefined{Theorem}{\newtheorem{Theorem}{Theorem}}{}
\@ifundefined{Proposition}{\newtheorem{Proposition}[Theorem]{Proposition}}{}
\makeatother
\title[The Fefferman Metric for Twistor CR Manifolds and Conformal ]{The Fefferman Metric for Twistor CR Manifolds and Conformal Geodesics in Dimension Three}
\author{Taiji Marugame}
\date{Source: arXiv:2411.18961v4 (CC BY 4.0)}
\begin{document}
\maketitle

\noindent\emph{Source and licence.} The Fefferman Metric for Twistor CR Manifolds and Conformal Geodesics in Dimension Three. Version arXiv:2411.18961v4 (CC BY 4.0), distributed under the Creative Commons Attribution 4.0 International licence, as recorded in the arXiv licence metadata for this exact version and shown on the abstract page https://arxiv.org/abs/2411.18961v4. Full rights notice: licenses/arxiv-2411.18961-CC-BY-4.0.txt.\par\smallskip

\noindent\textit{Editorial scope.} Counted unit: the Proposition whose source label is \texttt{vertical-variation} in the article (source0.tex lines 1175--1181, source label \texttt{vertical-variation}), delivered together with the article's own complete original proof of it (source0.tex lines 1182--1255, one \texttt{proof} environment of 74 lines). No mathematics is written, repaired or re-derived: statement and proof are the article's own text line for line. Required context retained uncounted: none. Every definition, display and label of those lines is retained unchanged. Omitted from this package and not redistributed: 1--1174, 1256--1561. Nothing inside a retained excerpt is omitted or altered: every retained body is delivered exactly as the article's lines are written, so no omission receipt of any kind accompanies this package. Citations: the retained excerpts carry no citation command at all, so no bibliography entry of the article is transcribed and no reference is redistributed. Reference adaptation: none was needed and none was made; every reference inside the delivered unit resolves inside it, so no unresolved reference or citation is produced. Printed slips kept literal: no printed word, symbol, hypothesis or number of the counted statement or of its proof was altered. Layout and class: the article's own class is \texttt{sigma} with packages including ifpdf, mathrsfs, xy; the wrapper is the lane's pinned \texttt{amsart} editorial wrapper at 10pt on A4, which restates the theorem-like environments the retained text needs and adds the article's own macro definitions that the retained text needs (none), with extra wrapper packages (bm, bbm, dsfont, xspace, cancel, mathtools). The delivered numbering is the unit's own. Assets: no retained span contains a figure environment, a graphics-inclusion command, a PDF-page inclusion, a table or a TikZ picture; no image file is copied into this package, the package declares no asset dependency and no image file is redistributed with it. Licence: version arXiv:2411.18961v4 of this article is distributed under the Creative Commons Attribution 4.0 International licence, as recorded in the arXiv licence metadata for this exact version and shown on the abstract page https://arxiv.org/abs/2411.18961v4. A full rights notice is delivered at licenses/arxiv-2411.18961-CC-BY-4.0.txt. Characters: every retained excerpt is pure ASCII, so the delivered engine, XeLaTeX with its default Latin Modern text font, reports no missing character.
\begin{Proposition}\label{vertical-variation}
Let $\gamma(t)=(x(t), B(t))$, $a\le t\le b$, be a curve on $\mathcal{P}$ which projects to $\tilde x(t)=(x(t), \dot x(t)/|\dot x(t)|)$. Let $\gamma_\epsilon(t)$ be a variation of $\gamma(t)$ fixing the endpoints such that the variation vector field $(\partial/\partial \epsilon)|_{\epsilon=0}\gamma_\epsilon(t)$ is vertical with respect to the fiberation $\mathcal{P}\to \Sigma$. Then, we~have
\[
\frac{\mathrm{d}}{\mathrm{d} \epsilon}\Big|_{\epsilon=0}
\int _a^b \frac{G^{\mathrm{F}}(\dot \gamma_\epsilon(t), \dot \gamma_\epsilon(t))}{G^{\mathrm{F}}(K, \dot \gamma_\epsilon(t))}\,{\rm d}t=0.
\]
\end{Proposition}

\begin{proof}
By the assumption on $\gamma(t)$, we have
\begin{equation}\label{phi-gamma}
\phi(\dot\gamma(t))=\begin{pmatrix} 0 \\ 0 \\ |\dot x|\end{pmatrix}.
\end{equation}

We may assume that the variation is of the form
\[
\gamma_\epsilon(t)=R_{{\rm e}^{\epsilon X(t)}}\gamma(t)=\gamma(t)\cdot {\rm e}^{\epsilon X(t)},
\]
where
\[
X(t)=\begin{pmatrix}
\hphantom{-}0 & \hphantom{-}X^3(t) & -X^2(t) \\
-X^3(t) & \hphantom{-}0 & \hphantom{-}X^1(t) \\
\hphantom{-}X^2(t) & -X^1(t) & \hphantom{-}0
\end{pmatrix}
\]
is an $\mathfrak{so}(3)$-valued function with $X(a)=X(b)=O$.

For any $X\in\mathfrak{so}(3)$, we denote by $X^\dagger$ the vertical vector field on $\mathcal{P}$ generated by $X$:
\[
X^\dagger_u:=\frac{\partial}{\partial s}\biggl|_{s=0}R_{{\rm e}^{sX}} u, \qquad u\in \mathcal{P}.
\]
Then, we have
\[
(R_{A})_* X^\dagger=\bigl({\rm Ad}\bigl(A^{-1}\bigr)X\bigr)^\dagger
\]
for any $A\in{\rm SO}(3)$ and $X\in\mathfrak{so}(3)$.

The velocity of the curve $\gamma_\epsilon(t)$ is given by
\[
\dot\gamma_\epsilon(t)=(R_{{\rm e}^{\epsilon X(t)}})_* \dot\gamma(t)+\epsilon\dot X(t)^{\dagger}.
\]
Then, by using the right-invariance of $G^{\rm F}$ and the fact that the fibers of $\mathcal{P}$ are totally null, we~compute as
\begin{align*}
G^{\rm F}(\dot\gamma_\epsilon, \dot\gamma_\epsilon)
={}&G^{\rm F} ((R_{{\rm e}^{\epsilon X}})_* \dot\gamma, (R_{{\rm e}^{\epsilon X}})_* \dot\gamma )
 2\epsilon G^{\rm F}\bigl((R_{{\rm e}^{\epsilon X}})_* \dot\gamma,
\dot X^{\dagger}\bigr)+\epsilon^2G^{\rm F}\bigl(\dot X^\dagger, \dot X^\dagger\bigr) \\
={}&G^{\rm F}(\dot\gamma, \dot\gamma)
+2\epsilon G^{\rm F}\bigl(\dot\gamma, \bigl({\rm Ad}\bigl({\rm e}^{\epsilon X}\bigr)\dot X\bigr)^\dagger\bigr).
\end{align*}
Hence, using \eqref{phi-gamma}, we have
\begin{equation*}%\label{var-numerator}
\frac{\partial}{\partial \epsilon}\bigg|_{\epsilon=0}G^{\rm F}(\dot\gamma_\epsilon, \dot\gamma_\epsilon)
=2G^{\rm F}\bigl(\dot\gamma, \dot X^\dagger\bigr)
=2|\dot x|\omega^3\bigl(\dot X^\dagger\bigr)=2|\dot x|\dot X^3.
\end{equation*}
On the other hand, since $K=Y^\dagger$ the denominator of the integrand is computed as
\begin{align*}
G^{\rm F}(K, \dot\gamma_\epsilon)
=G^{\rm F}\bigl(Y^\dagger, (R_{{\rm e}^{\epsilon X}})_* \dot\gamma+\epsilon\dot X^{\dagger}\bigr)
=G^{\rm F}\bigl(\bigl({\rm Ad}\bigl({\rm e}^{\epsilon X}\bigr) Y\bigr)^\dagger, \dot\gamma\bigr)
=\omega^3\bigl(\bigl({\rm Ad}\bigl({\rm e}^{\epsilon X}\bigr) Y\bigr)^\dagger\bigr),
\end{align*}
and hence
\[
\frac{\partial}{\partial \epsilon}\bigg|_{\epsilon=0}G^{\rm F}(K, \dot\gamma_\epsilon)
=\omega^3\bigl([X, Y]^\dagger\bigr) =\omega^3\left(\begin{pmatrix}
0 & 0 & -X^1 \\
0 & 0 & -X^2 \\
X^1 & X^2 & \hphantom{-}0
\end{pmatrix}^{\dagger}\right)=0.
\]
Thus, noting that $G^{F}(K, \dot\gamma)=|\dot x|\omega^3\bigl(Y^\dagger\bigr)=|\dot x|$, we have
\[
\frac{\mathrm{d}}{\mathrm{d} \epsilon}\bigg|_{\epsilon=0}
\int _a^b \frac{G^{\rm F}(\dot \gamma_\epsilon(t), \dot \gamma_\epsilon(t))}{G^{\rm F}(K, \dot \gamma_\epsilon(t))}\,{\rm d}t=
2\int_a^b \dot X^3(t) \, {\rm d}t=2\bigl(X^3(b)-X^3(a)\bigr)=0.
\tag*{\qed}
\]
\renewcommand{\qed}{}
\end{proof}

\end{document}
